\chapter{Challenges Faced and How They Were Addressed}\label{ch:challenges}

This chapter describes the main problems met in Phase I and how each was solved. Every decision is recorded with its
evidence in the project's decision log, including the ideas that did not work.


\section{Data}

\begin{xltabular}{\textwidth}{@{}P{32mm}LL@{}}
\caption{Challenges with the data.}\label{tab:ch-data}\\
\toprule
\textbf{Challenge} & \textbf{What happened} & \textbf{How we addressed it}\\
\midrule
\endfirsthead
\toprule
\textbf{Challenge} & \textbf{What happened} & \textbf{How we addressed it}\\
\midrule
\endhead
\bottomrule
\endlastfoot
\small The data did not match its documentation & \small The paper says 250 Hz and 886 seizures; the files say 256 Hz
and 883. Centre and age labels, needed for stratified splits, are not in the public release; seizure-visibility labels
are empty. & \small A full audit of all 11,009 files before any modelling. Every window and filter is defined in
seconds and designed for 256 Hz; folds are stratified by seizure count; analyses the data cannot support were dropped
and reported.\\
\small 126 GB of data on a machine without the disk or a GPU & \small The dataset could not be stored where the code
runs. & \small The audit and the feature extraction stream each file from the public S3 bucket, process it and delete
it; only the 3.5 GB of integer features are kept.\\
\small Silent download truncation & \small Comparing integer and floating-point features exposed one ECG file cut at
51,855 of 71,835 s. A new verifier then found 9 recordings silently truncated by incomplete network reads (one EEG file
had 12 of 16,363 windows). & \small Downloads are now checked against the expected size and every cached recording
against the audit durations. All 9 were re-extracted, and the 3 model runs trained on the bad cache were repeated.\\
\small A recording with an ECG lead but no detectable heartbeat & \small The feature extractor crashed on an empty list
of beats (3,250 s of data). & \small The recording is kept, its ECG features are set to 0 and a quality bit is
raised.\\
\small The official artefact rule does not suit behind-the-ear EEG & \small The 13\,\textmu V RMS floor flags 69\,\% of
ordinary windows on one channel, because behind-the-ear amplitudes are close to it. & \small For our models the
artefact flags are input bits, and whether to gate alarms on them is chosen on the training folds.\\
\end{xltabular}

\section{Reproducing the Published Baseline}

\begin{xltabular}{\textwidth}{@{}P{32mm}LL@{}}
\caption{Challenges in reproducing the published baseline.}\label{tab:ch-repro}\\
\toprule
\textbf{Challenge} & \textbf{What happened} & \textbf{How we addressed it}\\
\midrule
\endfirsthead
\toprule
\textbf{Challenge} & \textbf{What happened} & \textbf{How we addressed it}\\
\midrule
\endhead
\bottomrule
\endlastfoot
\small ChronoNet did not reproduce & \small Our run detected 74.5\,\% of seizures against 84.2\,\% published, with 22.6
instead of 100.5 false alarms per hour. & \small We ported the official code literally (identical outputs on 900 random
test cases) and then checked the metric definitions: the published AUPR matches only when averaged per recording
(0.109 against 0.104), and the published false-alarm rate matches neither definition in the code. The gap is reported
as measured.\\
\small The released code selects its model on the test patients & \small Its training loop keeps the epoch with the
best score on the same patients it is then evaluated on. & \small We reproduced it as published and also trained a
clean version (last epoch, no test information), which is the honest reference for the final comparison.\\
\small Defects in the released evaluation code & \small True positives can be counted more than once, and the artefact
gate (\texttt{a and b} on two Python lists) silently uses only one channel. & \small Reproduced as written for the
comparison with the paper; corrected numbers reported in a separate column.\\
\small The SVM baseline cannot be rebuilt yet & \small The dataset paper defers its 42 features to Vandecasteele et al.\
\cite{vandecasteele2020bte}, whose full text was not available to us. & \small Still open: we need that paper's feature
table.\\
\end{xltabular}

\section{Hardware Arithmetic}

\begin{xltabular}{\textwidth}{@{}P{32mm}LL@{}}
\caption{Challenges with the integer (hardware) arithmetic.}\label{tab:ch-hw}\\
\toprule
\textbf{Challenge} & \textbf{What happened} & \textbf{How we addressed it}\\
\midrule
\endfirsthead
\toprule
\textbf{Challenge} & \textbf{What happened} & \textbf{How we addressed it}\\
\midrule
\endhead
\bottomrule
\endlastfoot
\small Rounding noise in integer filters & \small With plain rounding, the 0.5 Hz high-pass filter had an error of 140
LSB standard deviation, larger than the signal itself (61 LSB). & \small Second-order error feedback of the discarded
bits reduced the error to 0.7 LSB.\\
\small Quantised coefficients moved the filter zeros & \small Rounding each coefficient independently moved the
low-pass zeros away from the Nyquist frequency (stop-band floor $-55$ dB). & \small Structure-preserving quantisation,
$b = b_0\,(1, \pm 2, 1)$, with coefficients as sums of powers of two.\\
\small Heart rate without a divider & \small Heart rate needs $60/\mathrm{RR}$. & \small We use
$\log_2(\mathrm{RR})$ from integer R-peak timestamps: no division and no multiplication.\\
\small Proving that the integer pipeline costs nothing & \small Hardware arithmetic could lose accuracy silently. &
\small A floating-point twin of every feature: median correlation 0.9965 and equal LightGBM AUROC (0.736 against
0.735).\\
\small TreeLUT scores were not comparable across folds & \small Each fold's integer score had its own scale, so pooled
thresholds were meaningless; 3-bit leaves could not represent the model. & \small Scores converted to a common unit;
leaf width chosen on the folds (5 bits).\\
\small Verilog disagreed with the software & \small Thresholds outside a feature's bit width (for example 79 on a 2-bit
presence feature) produced wrong comparator logic. & \small The Verilog generator now emits constant bits for such
thresholds; the simulation agrees with the software on 100\,\% of the test vectors.\\
\end{xltabular}

\section{Computing and Infrastructure}

\begin{xltabular}{\textwidth}{@{}P{32mm}LL@{}}
\caption{Challenges with computing resources.}\label{tab:ch-infra}\\
\toprule
\textbf{Challenge} & \textbf{What happened} & \textbf{How we addressed it}\\
\midrule
\endfirsthead
\toprule
\textbf{Challenge} & \textbf{What happened} & \textbf{How we addressed it}\\
\midrule
\endhead
\bottomrule
\endlastfoot
\small Limited GPU resources & \small Our GPU is a laptop RTX 3060 with 6 GB of memory and 7.4 GB of system memory,
and the reference DWN library runs only on CUDA. & \small We re-implemented the DWN in PyTorch, identical to the
library's CUDA kernel (unit-tested), so that it also runs on CPUs. CPU jobs run on a 4-core server; large networks and
raw-EEG models run on the GPU.\\
\small Long runs interrupted & \small Machine restarts stopped long runs; the logic-gate network lost 3 of its 5
folds once. & \small Every run resumes at fold and epoch level, and a queue skips finished runs; a resumed run was
checked to give identical metrics.\\
\small Out-of-memory failures & \small The stateful network's sequence sampler held 24 million windows in memory
(13.9 GB of a 15 GB machine). & \small A new sampler reads each recording once and stores bits as bytes: the same
sequences with a peak of 3.1 GB.\\
\small Two machines, one code base & \small The CPU server and the GPU laptop must not overwrite each other's work. &
\small One version-controlled code base, written run instructions for each machine, and automatic tests on every
change.\\
\end{xltabular}

\section{Machine Learning}

\begin{xltabular}{\textwidth}{@{}P{32mm}LL@{}}
\caption{Challenges in the machine-learning design.}\label{tab:ch-ml}\\
\toprule
\textbf{Challenge} & \textbf{What happened} & \textbf{How we addressed it}\\
\midrule
\endfirsthead
\toprule
\textbf{Challenge} & \textbf{What happened} & \textbf{How we addressed it}\\
\midrule
\endhead
\bottomrule
\endlastfoot
\small Seizures are 0.123\,\% of the time & \small Accuracy is meaningless, and training on all windows teaches the
model never to fire. & \small Event-based scoring (SzCORE) per patient; training on all seizure windows plus a fixed
ratio of background windows; operating points chosen for a false-alarm budget.\\
\small Results must hold on unseen patients & \small Mixing a patient's windows between training and validation
inflates results. & \small Patient-wise five-fold cross-validation; feature selection and thresholds fitted inside each
training fold only; the test patients locked behind a flag, with every access logged.\\
\small Early results were weaker than expected & \small The first LUT network, with mutual-information features,
reached an event-F1 of 0.148 with 7.0 false alarms per day; the interpretability study showed that it ignored the
heart-rate and muscle features. & \small Feature selection by LightGBM importance: event-F1 0.208 and 2.3 false alarms
per day at the same size.\\
\small A confounded comparison & \small The stateful network had received all 117 features instead of the selected 40,
so its apparent gain mixed two changes. & \small Found by checking the run's records; all models now share one training
function, and the run was repeated.\\
\small Stacking past windows picked duplicate features & \small Ranking the 585 stacked columns kept lagged copies of
only 11--13 distinct features. & \small Rank on the current window, then keep every lag of the selected features.\\
\end{xltabular}
